\documentclass[11pt,a4paper]{article}

\usepackage[margin=2.5cm]{geometry}
\usepackage{amsmath,amssymb,amsthm}
\usepackage{bm}
\usepackage{booktabs}
\usepackage{xcolor}
\usepackage{hyperref}
\usepackage{enumitem}

\newtheorem{remark}{Remark}

\newcommand{\R}{\mathbb{R}}
\newcommand{\dom}{\Omega}
\newcommand{\pore}[1]{\Omega_{#1}}
\newcommand{\iface}{\gamma}
\newcommand{\vel}{\bm{u}}
\newcommand{\pres}{p}
\newcommand{\stress}{\bm{\sigma}}
\newcommand{\normal}{\bm{n}}
\newcommand{\trac}{\bm{\lambda}}
\newcommand{\DtN}{\mathbf{G}}
\newcommand{\Schur}{\mathbf{S}}
\newcommand{\one}{\mathbf{1}}
\newcommand{\zero}{\bm{0}}

\newcommand{\modeset}{\mathcal{M}}

\hypersetup{colorlinks=true,linkcolor=blue,urlcolor=blue,citecolor=blue}

\begin{document}

\title{\textbf{Affine-DDPNM: Methodology}}
\author{}
\date{\today}
\maketitle

We extend the domain-decomposition pore-network method (DD-PNM) of
\cite{sun2025ddpnm} by enriching the interface traction space from a single
constant-normal mode to a nine-mode affine representation.  The method,
termed \textbf{Affine-DDPNM}, retains the offline--online workflow and
non-overlapping subdomain partition of the original DD-PNM.  The only
change is the interface Ansatz: on each internal face, the traction is
now allowed to vary as an affine function of the in-plane coordinates, in
both the normal and the two tangential directions.

% ======================================================================
\section{Model problem}
% ======================================================================

We consider steady, incompressible Stokes flow in a three-dimensional
pore domain $\dom \subset \R^3$.  Let $\Gamma_{\text{wall}}$ be the
solid--fluid interface.  A pressure drop $p_{\text{in}} - p_{\text{out}}
> 0$ is imposed between the inlet and outlet boundaries, while all
remaining external boundaries are treated as no-slip walls.  Collecting
all homogeneous Dirichlet segments in $\Gamma_{\text{dir}}$, the governing
equations are
\begin{align}
  -\nu\Delta\vel + \nabla\pres &= \bm{0},                && \text{in } \dom,\\
  \nabla\cdot\vel &= 0,                                    && \text{in } \dom,\\
  \vel &= \bm{0},                                          && \text{on } \Gamma_{\text{dir}},\\
  \normal\cdot\stress(\vel,\pres)\normal &= -p_b,\quad
    (I - \normal\otimes\normal)\,\stress(\vel,\pres)\normal = \bm{0},
                                                          && \text{on } \Gamma_{\text{in}}\cup\Gamma_{\text{out}},
\end{align}
where $\vel$ is velocity, $\pres$ is pressure, $\nu$ is the kinematic
viscosity, $\normal$ is the outward unit normal, and $p_b = p_{\text{in}}$
on $\Gamma_{\text{in}}$, $p_b = p_{\text{out}}$ on $\Gamma_{\text{out}}$.
The \emph{pseudostress} tensor is
\[
  \stress(\vel,\pres) := -\pres\mathbf{I} + \nu\nabla\vel,
\]
which for incompressible flow differs from the standard Cauchy stress
$-\pres\mathbf{I} + 2\mu\bm{\varepsilon}(\vel)$ in the boundary traction,
but the two coincide in the interior momentum equation.
The zero tangential traction condition $(I - \normal\otimes\normal)\,
\stress(\vel,\pres)\normal = \bm{0}$ is equivalent to
$\bm{t}_{\nu}\cdot\stress(\vel,\pres)\normal = 0$ for two orthonormal
tangents $\bm{t}_1,\bm{t}_2$ spanning the boundary plane.

% ======================================================================
\section{Domain partition}
% ======================================================================

The pore domain $\dom$ is partitioned into $N$ non-overlapping Lipschitz
subdomains $\{\pore{i}\}_{i=1}^N$, each corresponding to one pore body.
For any pair $(i,j)$ of adjacent pores, the shared interface is
\[
  \iface_{ij} = \partial\pore{i} \cap \partial\pore{j},
\]
and we index the $m$ internal interfaces as $\{\iface_k\}_{k=1}^m$.
Interfaces are assumed \textbf{planar} throughout the theoretical
development; the numerical watershed experiments relax this and are
discussed separately in Remark~\ref{rem:planar}.

For each subdomain $\pore{i}$ we list its $m_i$ ports as
$\{\iface_{i,1},\dots,\iface_{i,m_i}\}$, partitioned into an
unknown-traction set $\mathcal{U}_i$ (internal interfaces) and a
known-traction set $\mathcal{K}_i$ (inlet/outlet boundaries).  The
local-to-global index maps and the Boolean restriction matrices
$R_i^u$, $R_i^k$ follow the convention of \cite[Section~2.4]{sun2025ddpnm}.

The Voronoi partition places each interface as the plane perpendicular to
and bisecting the line connecting the two neighbouring sphere centres;
for unequal radii this plane differs from the exact clearance saddle
surface.  The geometric error introduced by this approximation is
quantified in the numerical companion.

% ======================================================================
\section{Reference finite-element discretization}
% ======================================================================

Within each subdomain we employ the inf--sup stable Taylor--Hood pair
$\mathbf{V}_h = [P_2]^3 \cap \mathbf{V}$,
$Q_h = P_1 \cap Q$ on a body-fitted tetrahedral mesh, where
\[
  \mathbf{V} = \{\bm{v} \in [H^1(\dom)]^3 : \bm{v} = \bm{0}
    \text{ on } \Gamma_{\text{dir}}\},\qquad
  Q = L^2(\dom).
\]
With bases $\{\bm{\phi}_j\}_{j=1}^{n_u}$ and $\{\psi_l\}_{l=1}^{n_p}$,
the discrete system on $\pore{i}$ under Neumann boundary data reads
\begin{equation}\label{eq:local-fe}
  \begin{bmatrix} A_i & B_i^\top \\ B_i & 0 \end{bmatrix}
  \begin{bmatrix} \mathbf{U}_i \\ \mathbf{P}_i \end{bmatrix}
  = \begin{bmatrix} \sum_{r=1}^{m_i} \bm{f}_i^{(r)} \\ \bm{0} \end{bmatrix},
\end{equation}
where $A_i$ is the velocity-stiffness matrix (SPD),
$B_i$ is the divergence matrix, and $\bm{f}_i^{(r)}$ is the Neumann load
vector assembled from the traction prescribed on port $r$.
No pressure stabilisation is used ($\varepsilon = 0$).

% ======================================================================
\section{Interface traction Ansatz}
% ======================================================================

\subsection{Classic DD-PNM (one mode per interface)}

In the classic DD-PNM \cite{sun2025ddpnm}, the traction on each internal
interface $\gamma$ is uniform and purely normal:
\begin{equation}\label{eq:classic-traction}
  -\stress(\vel,\pres)\normal_\gamma = p_\gamma\,\normal_\gamma,
\end{equation}
where $p_\gamma \in \R$ is a single scalar per interface.  Equivalently,
$\mathcal{T}_\gamma^{\text{P0}} = \operatorname{span}\{\normal_\gamma\}$.

\subsection{Affine-DDPNM (nine modes per interface)}
\label{sec:affine-ansatz}

For each \textbf{planar} interface $\gamma$, we introduce a local
orthonormal frame $\{\normal_\gamma, \bm{t}_{\gamma,1},
\bm{t}_{\gamma,2}\}$ with origin at the interface centroid
$\bm{c}_\gamma$.  The dimensionless in-plane coordinates are
\begin{equation}\label{eq:local-coords}
  s(\bm{x}) = \frac{(\bm{x} - \bm{c}_\gamma) \cdot
    \bm{t}_{\gamma,1}}{a_\gamma},\qquad
  t(\bm{x}) = \frac{(\bm{x} - \bm{c}_\gamma) \cdot
    \bm{t}_{\gamma,2}}{b_\gamma},
\end{equation}
where $a_\gamma, b_\gamma$ are the half-spans of $\gamma$ along the two
tangent directions, ensuring $s,t = O(1)$ under shape-regularity.

For a pore $\pore{i}$ adjacent to $\gamma$, define the \emph{signed}
direction vectors
\begin{equation}\label{eq:signed-dir}
  \bm{d}_{i,\gamma,\normal} = \normal_{i,\gamma},\qquad
  \bm{d}_{i,\gamma,\bm{t}_\nu} = \sigma_{i,\gamma}\,\bm{t}_{\gamma,\nu}
  \;\;(\nu=1,2),
\end{equation}
where $\sigma_{i,\gamma}=+1$ if $\pore{i}$ is the first member of the
interface pair and $-1$ otherwise.  These satisfy
$\bm{d}_{j,\gamma,\beta} = -\bm{d}_{i,\gamma,\beta}$, guaranteeing
that continuity conditions written below are sign-consistent.

The applied traction on $\gamma$ (as seen from $\pore{i}$) is expanded as
\begin{equation}\label{eq:affine-traction}
  -\stress(\vel,\pres)\,\normal_{i,\gamma}
  = \sum_{(\alpha,\beta)\in\modeset}
    \lambda_{\gamma,\alpha,\beta}\;
    \varphi_\alpha(s,t)\,\bm{d}_{i,\gamma,\beta},
\end{equation}
where the \textbf{mode index set} is
$\modeset = \{0,s,t\} \times \{\normal,\bm{t}_1,\bm{t}_2\}$
with scalar shape functions
$\varphi_0(s,t) = 1$, $\varphi_s(s,t) = s$, $\varphi_t(s,t) = t$.
The minus sign on the left follows the classic convention
$-\sigma n = p\,n$ and ensures that a positive $\lambda_{0,\normal}$
corresponds to a pressure acting into the pore.

Each interface carries \textbf{nine} scalar coefficients
$\trac_\gamma = (\lambda_{\gamma,0,\normal},\,
\lambda_{\gamma,0,\bm{t}_1},\,\lambda_{\gamma,0,\bm{t}_2},\,
\lambda_{\gamma,s,\normal},\,\dots,\,
\lambda_{\gamma,t,\bm{t}_2}) \in \R^9$
as the unknown degrees of freedom.
The affine traction space is
\[
  \mathcal{T}_\gamma^{\text{Affine}} = \operatorname{span}\{\,
    \varphi_\alpha(s,t)\,\bm{d}_{i,\gamma,\beta}
    \mid (\alpha,\beta) \in \modeset \,\}.
\]

\begin{remark}[Behaviour of the nine modes]\label{rem:modes}
  The constant normal mode $(\alpha,\beta)=(0,\normal)$ corresponds to
  a uniform normal pressure on the interface, matching the classic
  single-mode Ansatz.  The linear normal modes $(s,\normal)$ and
  $(t,\normal)$ capture first-order spatial variation of the normal
  traction.  The six tangential modes $(\alpha,\bm{t}_\nu)$ allow
  non-zero shear tractions at the interface, representing cross-flow
  and viscous momentum exchange that the purely normal classic Ansatz
  cannot express.
\end{remark}

\begin{remark}[Planar interface requirement]\label{rem:planar}
  The theory assumes planar interfaces with a fixed $\normal_\gamma$.
  Watershed partitions produce bent, faceted interfaces where the
  mesh-level facet normal $\normal(x)$ varies.  In the code,
  watershed interfaces use the per-facet $\normal(x)$ for the normal
  direction but a fixed averaged frame for the tangents; this
  deviates from the nine-mode space defined here and is one factor
  behind the degraded Affine accuracy on watershed geometries.
  All theoretical statements below apply to planar interfaces.
\end{remark}

% ======================================================================
\section{Local compliance map}
% ======================================================================

\subsection{Primitive response loads}

For each subdomain $\pore{i}$ and each mode $(\alpha,\beta) \in \modeset$
on an internal interface $\iface_{i,r}$, the Neumann load vector is
\begin{equation}\label{eq:unit-load}
  \bm{f}_{i,r}^{(\alpha,\beta)} \in \R^{n_{i,u}},\qquad
  \bigl[\bm{f}_{i,r}^{(\alpha,\beta)}\bigr]_j
  = \int_{\iface_{i,r}} \varphi_\alpha(s,t)\;
    \bm{d}_{i,\gamma,\beta} \cdot \bm{\phi}_j^{(i)} \;\mathrm{d}S.
\end{equation}
With the sign convention~\eqref{eq:affine-traction}
$-\sigma n = \sum \lambda\varphi\bm{d}$, a positive coefficient
$\lambda$ produces a traction acting \emph{into} the pore; the load
vector thus carries no leading minus sign---it directly provides the
right-hand side of~\eqref{eq:local-fe} for the applied boundary force.
For boundary ports only the uniform-normal mode
$(\alpha,\beta)=(0,\normal)$ is used, with a \emph{prescribed}
coefficient equal to the inlet/outlet pressure.

Solving the local Stokes system for each unit load yields the velocity
response $\mathbf{U}_{i,r}^{(\alpha,\beta)}$ and pressure response
$\mathbf{X}_{i,r}^{(\alpha,\beta)}$:
\begin{equation}\label{eq:response}
  \begin{bmatrix} A_i & B_i^\top \\ B_i & 0 \end{bmatrix}
  \begin{bmatrix} \mathbf{U}_{i,r}^{(\alpha,\beta)} \\
                   \mathbf{X}_{i,r}^{(\alpha,\beta)} \end{bmatrix}
  = \begin{bmatrix} \bm{f}_{i,r}^{(\alpha,\beta)} \\ \bm{0} \end{bmatrix}.
\end{equation}

\subsection{Compliance (Neumann-to-Dirichlet) map}

Let $n_i = 9m_i^u + m_i^k$ be the total number of local traction
coefficients (nine per internal interface, one per boundary port;
the $m_i^k$ boundary entries are \emph{prescribed}, not unknown).
Collect them in $\trac_i \in \R^{n_i}$.

By linearity of~\eqref{eq:local-fe}, the discrete velocity is
\begin{equation}\label{eq:velocity-linear}
  \mathbf{U}_i
  = \sum_{\gamma\in\mathcal{U}_i}\sum_{\alpha,\beta}
    \lambda_{i,\gamma,\alpha,\beta}\,
    \mathbf{U}_{i,\gamma}^{(\alpha,\beta)}
  \;+\; \sum_{\gamma\in\mathcal{K}_i}
    \lambda_{i,\gamma}^k\,
    \mathbf{U}_{i,\gamma}^{(0,\normal)}.
\end{equation}

For each mode $(\alpha,\beta)$ on interface $\iface_{i,r}$, the
\emph{weighted velocity moment} is
\begin{equation}\label{eq:moment}
  M_{i,r}^{(\alpha,\beta)}
  = \int_{\iface_{i,r}} \varphi_\alpha(s,t)\;
    \vel_{i,h} \cdot \bm{d}_{i,\gamma,\beta} \;\mathrm{d}S.
\end{equation}
Only $M_{i,r}^{(0,\normal)} = \int_\gamma \vel\cdot\normal\,\mathrm{d}S$
is the volumetric flow rate.  The remaining eight quantities are
weighted velocity moments: $M^{(s,\normal)},M^{(t,\normal)}$ are
linear normal flux moments, and
$M^{(\alpha,\bm{t}_1)},M^{(\alpha,\bm{t}_2)}$
are tangential velocity moments.  Collectively we call them the
\textbf{nine weighted velocity moments}.

The local compliance (Neumann-to-Dirichlet) map
$\DtN_i \in \R^{n_i \times n_i}$ has entries
\begin{equation}\label{eq:dtn}
  \DtN_i\bigl[(\gamma,\alpha,\beta),\,(\gamma',\alpha',\beta')\bigr]
  = \bigl[\bm{f}_{i,\gamma}^{(\alpha,\beta)}\bigr]^\top\,
    \mathbf{U}_{i,\gamma'}^{(\alpha',\beta')}.
\end{equation}
For boundary ports, only the $(\alpha,\beta)=(0,\normal)$ entry is stored.
The map is symmetric positive semi-definite; the energy identity and
kernel characterisation are given in the companion mathematical
properties document.

$\DtN_i$ is block-partitioned into $\DtN_i^{uu}$ (unknown--unknown) and
$\DtN_i^{uk}$ (unknown--known) as in the classic case.

% ======================================================================
\section{Global Schur complement assembly}
% ======================================================================

Let $\trac \in \R^{9m}$ collect all unknown interface coefficients and
$\trac^k \in \R^{m^k}$ the prescribed boundary tractions.
Enforcing continuity of \textbf{all nine weighted velocity moments} on
every internal interface yields the global Schur system
\begin{equation}\label{eq:schur}
  \Schur\,\trac = \bm{F},\qquad
  \Schur = \sum_{i=1}^N (R_i^u)^\top \DtN_i^{uu} R_i^u,\qquad
  \bm{F} = -\sum_{i=1}^N (R_i^u)^\top \DtN_i^{uk} R_i^k \trac^k,
\end{equation}
where $\Schur \in \R^{9m \times 9m}$.  The continuity condition
$M_{i,r}^{(\alpha,\beta)} + M_{j,r}^{(\alpha,\beta)} = 0$ follows from
$\bm{d}_{j,\gamma,\beta} = -\bm{d}_{i,\gamma,\beta}$ and is imposed
for all nine modes, not only the constant normal component.

\begin{remark}[SPD of the global system]\label{rem:spd}
  $\Schur$ is symmetric by construction.  Positive definiteness requires:
  \begin{enumerate}[nosep,leftmargin=*]
    \item Every connected component of the pore adjacency graph contains
      at least one pore with a prescribed-traction (inlet/outlet) port
      (the \emph{anchoring condition});
    \item Each pore has a positive-measure no-slip wall segment, fixing
      rigid-body modes;
    \item The nine primitive loads per interface are linearly independent
      on each face (verified numerically for all tested geometries);
    \item Interface coordinate frames and side signs are consistent across
      adjacent pores.
  \end{enumerate}
  Under these conditions, the Schur system is SPD and uniquely solvable.
  In practice, watershed partitions may produce \emph{floating basins}
  (pores with no solid-wall facet); these are merged into their
  neighbours before the Schur assembly to enforce the anchoring condition.
\end{remark}

% ======================================================================
\section{Solution algorithm}
\label{sec:algorithm}

The Affine-DDPNM workflow is summarised in Algorithm~\ref{alg:affine}.
The differences from the classic algorithm are: (i)~nine primitive loads
per internal interface instead of one, and (ii)~the $9\times$ increase
in the size of the local compliance blocks and the global Schur system.
The offline--online decomposition and the embarrassingly parallel
structure are preserved unchanged.

\bigskip
\noindent\framebox{%
\begin{minipage}{0.95\textwidth}
\textbf{Algorithm~\ref{alg:affine}:} Affine-DDPNM solver
(assemble -- couple -- recover)
\label{alg:affine}

\smallskip
\noindent\textbf{Input:} subdomain meshes, FE pairs $(V_{i,h}, Q_{i,h})$,
  prescribed tractions $\trac^k$.

\noindent\textbf{Output:} unknown interface coefficients $\trac$,
  weighted velocity moments, and (optionally) local fields.

\smallskip
\begin{enumerate}[nosep,leftmargin=*]
  \item[\bfseries 1.] \textbf{Offline library} --- For each pore body
    $i = 1,\dots,N$ (in parallel):
    \begin{enumerate}[nosep,leftmargin=*]
      \item Assemble local operators $A_i$, $B_i$.
      \item For each internal interface $\gamma \in \mathcal{U}_i$ and
        each mode $(\alpha,\beta) \in \modeset$:
        assemble unit load $\bm{f}_{i,\gamma}^{(\alpha,\beta)}$ via
        \eqref{eq:unit-load}, solve~\eqref{eq:response} for
        $(\mathbf{U}_{i,\gamma}^{(\alpha,\beta)},
          \mathbf{X}_{i,\gamma}^{(\alpha,\beta)})$.
      \item For each boundary port $\gamma \in \mathcal{K}_i$:
        assemble and solve the single uniform-normal load
        $(\alpha,\beta) = (0,\normal)$.
      \item Form local compliance block $\DtN_i$ via \eqref{eq:dtn}.
      \item Accumulate $S \mathrel{+}= (R_i^u)^\top \DtN_i^{uu} R_i^u$,
        $\bm{F} \mathrel{-}= (R_i^u)^\top \DtN_i^{uk} R_i^k \trac^k$.
    \end{enumerate}
  \item[\bfseries 2.] \textbf{Global solve} --- Solve
    $\Schur\,\trac = \bm{F}$ by sparse direct solver.
  \item[\bfseries 3.] \textbf{Reconstruction} (optional, in parallel):
    \begin{align*}
      \mathbf{U}_i &= \sum_{\gamma\in\mathcal{U}_i}\sum_{\alpha,\beta}
        \lambda_{i,\gamma,\alpha,\beta}\,
        \mathbf{U}_{i,\gamma}^{(\alpha,\beta)}
        \;+\; \sum_{\gamma\in\mathcal{K}_i} \lambda_{i,\gamma}^k\,
        \mathbf{U}_{i,\gamma}^{(0,\normal)},\\[2pt]
      \mathbf{P}_i &= \sum_{\gamma\in\mathcal{U}_i}\sum_{\alpha,\beta}
        \lambda_{i,\gamma,\alpha,\beta}\,
        \mathbf{X}_{i,\gamma}^{(\alpha,\beta)}
        \;+\; \sum_{\gamma\in\mathcal{K}_i} \lambda_{i,\gamma}^k\,
        \mathbf{X}_{i,\gamma}^{(0,\normal)}.
    \end{align*}
\end{enumerate}
\end{minipage}}
\bigskip

% ======================================================================
\section{Relation to the classic DD-PNM}
\label{sec:reduction}
% ======================================================================

Let $\mathcal{P}_{\text{P0}} : \R^m \to \R^{9m}$ be the embedding that
maps the $m$ classic constant-normal coefficients to the $9m$ affine
space by placing them in the $(0,\normal)$ slots and setting all other
components to zero.  The classic traction space
$\operatorname{Range}(\mathcal{P}_{\text{P0}})$ is a subspace of the
affine traction space.

The classic DD-PNM is recovered by \textbf{restricting both the trial
and test spaces} to the constant-normal subspace:
\begin{equation}\label{eq:classic-recovery}
  \Schur_{\text{P0}} = \mathcal{P}_{\text{P0}}^\top \Schur\,
                       \mathcal{P}_{\text{P0}},\qquad
  \bm{F}_{\text{P0}} = \mathcal{P}_{\text{P0}}^\top \bm{F},
\end{equation}
yielding the classic $m\times m$ Schur system.  In general, the classic
solution $\bm{p} = \Schur_{\text{P0}}^{-1}\bm{F}_{\text{P0}}$ does
\emph{not} coincide with the $(0,\normal)$-components of the full
$9m$-system solution $\trac = \Schur^{-1}\bm{F}$, because the eight
additional test-space constraints alter the Galerkin projection.

The enrichment guarantee is that the affine space contains the classic
space as a subspace of the traction representation, and the Galerkin
optimality (proved in the companion mathematical properties document)
ensures that the affine error in the $\widehat{\Schur}$-energy norm is
never larger than the classic error.  The 73--90\% $L^2$ reductions
observed numerically are \emph{consistent} with this energy-norm
guarantee.

% ======================================================================
\section{Computational cost}
% ======================================================================

For a pore with $m_i$ ports, the local compliance matrix $\DtN_i$ has
size $n_i \times n_i$ where $n_i \approx 9m_i$ (versus $m_i$ in the
classic method).  Building $\DtN_i$ requires $n_i$ independent local
Stokes solves; this offline work is embarrassingly parallel.
The global Schur system $\Schur$ has size $9m \times 9m$ and is solved
by a sparse direct solver.

Offline cost grows approximately $9\times$ in the number of local solves.
Online cost is dominated by the sparse LU factorisation of $\Schur$,
which has at most $9 \times (\text{max coordination}) \times 9$
non-zero blocks per row.  For multi-boundary-condition studies the
offline library is built once and reused; the on-line solve then
achieves the speedups reported in the numerical companion
($100$--$10^4\times$ over monolithic FEM for Classic, $1.7$--$840\times$
for Affine).

\begin{thebibliography}{99}

\bibitem{sun2025ddpnm}
S.~Sun, Z.~Wang, L.~Zhang, J.~Zhao.
\newblock A domain decomposition approach to pore-network modeling of porous
  media flow.
\newblock \textit{arXiv:2510.13429v2}, 2026.

\end{thebibliography}

\end{document}
